\documentclass[10pt]{article}
\usepackage[margin=0.78in]{geometry}
\usepackage{amsmath,amssymb}
\usepackage[T1]{fontenc}
\usepackage[hidelinks]{hyperref}
\usepackage{titling}
\setlength{\droptitle}{-4em}

\title{Literature Status: A Complete Combinatorial Classification of Bounded Slice Spaces}
\author{fa\_banach\_001, lane 8}
\date{August 9, 2026}

\begin{document}
\maketitle

\paragraph{Status.}
\textbf{literature\_already\_answered}.  This is an exact answer explicitly identified by the supporting authors, not an agent-inferred implication.

\paragraph{Original source and question.}
Gaubert and Gunawardena, \emph{The Perron--Frobenius Theorem for Homogeneous, Monotone Functions}, arXiv:math/0105091, later Trans. Amer. Math. Soc. 356 (2004), 4931--4950.  In the subsection ``Slice spaces and recession functions,'' journal page 4937, after giving a sufficient recession-function condition, they ask:
\begin{quote}
It remains an open problem whether the boundedness of all slice spaces can be determined by combinatorial or graph-theoretic constructions as in Theorems 4 and 6.
\end{quote}
Here a slice space is
\[
  \mathcal S^\beta_\alpha(f)
  =\{x\in\mathbb R^n_{>0}:\alpha x\leq f(x)\leq \beta x\},
  \qquad \alpha,\beta>0,
\]
for a monotone positively homogeneous self-map $f$ of $\mathbb R^n_{>0}$, and boundedness is measured in Hilbert's projective metric.

\paragraph{Supporting answer.}
Akian, Gaubert, and Hochart, \emph{A game theory approach to the existence and uniqueness of nonlinear Perron--Frobenius eigenvectors}, arXiv:1812.09871.  Its introduction restates the source question verbatim in substance as Problem 1 [GG04, p.~4937] and says ``we ... solve Problem 1.''  The exact answer is Theorem 1.2: all slice spaces of $f$ are bounded in Hilbert's projective metric if and only if the two players have no disjoint dominions in the finite-state game $\Gamma_\infty(f)$.  The legal actions of this game are defined only from coordinatewise growth of the logarithmic conjugate $T=\log\circ f\circ\exp$ at $+\infty$ and $-\infty$, so this is a combinatorial criterion.  Section 4 further shows that the dominion condition is decidable by reachability in two directed hypergraphs.

\paragraph{Identification and scope.}
The supporting paper cites the original article, gives the exact original page and problem, and explicitly presents Theorem 1.2 as the full characterization solving it.  Thus the finite-dimensional slice-space problem is completely answered.  This identification does not settle the distinct question, also raised in the 2004 article, of extending these methods to general cones in Banach spaces.

\paragraph{Evidence searched.}
The run indexes contained no prior record for this paper or question.  An arXiv title-and-keyword search located arXiv:1812.09871; its abstract, Problem 1, Theorem 1.2, and Section 4 provide the decisive evidence.

\begin{thebibliography}{9}
\bibitem{GG04}
S. Gaubert and J. Gunawardena,
\emph{The Perron--Frobenius Theorem for Homogeneous, Monotone Functions},
arXiv:math/0105091; Trans. Amer. Math. Soc. 356 (2004), 4931--4950.

\bibitem{AGH}
M. Akian, S. Gaubert, and A. Hochart,
\emph{A game theory approach to the existence and uniqueness of nonlinear Perron--Frobenius eigenvectors},
arXiv:1812.09871.
\end{thebibliography}

\end{document}
